\chapter{RStudio, R, Variables, and Mathematical Operations}
\label{chap:r-basics}

\chapterguide
  {No prior R programming is assumed. The lab handout begins with installation and a one-line environment check.}
  {install and verify R/RStudio; distinguish Console and Script usage; create variables; use assignment operators; inspect object classes; perform arithmetic; and read simple user input.}

\section{R and RStudio setup}

The first handout separates the setup into two installations: install \term{R} first, then install \term{RStudio Desktop}. The environment is considered operational when RStudio opens and the Console evaluates a simple expression such as:

\begin{lstlisting}[style=dsr]
1 + 2
\end{lstlisting}

The expected result is \texttt{3}. This small test checks that RStudio can communicate with the installed R runtime.

\begin{keyidea}
R is the programming language/runtime used to execute the code. RStudio is the development environment shown in the lab handout. The lab sequence therefore installs R before RStudio.
\end{keyidea}

\subsection{Console versus Script}

The handout asks you to differentiate the R Console and an R Script. Their roles can be inferred from the exercises:

\begin{dstable}
\begin{tabularx}{\linewidth}{@{}p{3.0cm}X X@{}}
\toprule
\rowcolor{DSPaleBlue!92}
Area & What the lab uses it for & Example \\
\midrule
R Console & Immediate commands and quick experiments. A command is entered and evaluated directly. & \texttt{2+3}, \texttt{print("Hello World!")} \\
R Script & A saved sequence of R statements that can be rerun. & A file named \texttt{test.R} containing a short program. \\
\bottomrule
\end{tabularx}
\end{dstable}

\begin{workedexample}
The first saved script in the handout is:
\begin{lstlisting}[style=dsr]
# My first program in R Programming
myString <- "Hello, World!"
print(myString)
\end{lstlisting}
The line beginning with \texttt{\#} is a comment. The assignment stores a character value in \texttt{myString}, and \texttt{print()} displays it.
\end{workedexample}

\section{Packages}

The handout gives two package-installation paths: use the RStudio menus, or call \texttt{install.packages()} from the Console. A package is then made available in the current R session using \texttt{library()}.

\begin{lstlisting}[style=dsr]
install.packages("package_name")
library(package_name)
\end{lstlisting}

The installation step is distinct from loading: installation places the package on the system; \texttt{library()} attaches it for use in the current session.

\section{Variables and object names}

The lab creates three variables and inspects the current workspace:

\begin{lstlisting}[style=dsr]
var.1 = 5
var_1 = 7
x = 1
print(ls())
print(ls(pattern = "var"))
\end{lstlisting}

The first \texttt{ls()} lists the objects in the workspace. The pattern form restricts the names shown to those matching \texttt{"var"}. The exercise also demonstrates that names containing a period or underscore can be used.

\section{Assignment operators}

The supplied table groups assignment operators by direction:

\begin{dstable}
\begin{tabularx}{\linewidth}{@{}p{3.3cm}X@{}}
\toprule
\rowcolor{DSPaleBlue!92}
Operator & Handout description \\
\midrule
\texttt{<-}, \texttt{<<-}, \texttt{=} & Leftwards assignment \\
\texttt{->}, \texttt{->>} & Rightwards assignment \\
\bottomrule
\end{tabularx}
\end{dstable}

The examples deliberately vary assignment direction and value type:

\begin{lstlisting}[style=dsr]
Gender <- "Female"
height <<- 152
Weight = 81
3 -> f
23.5 ->> x
b <- seq(from = 0, to = 5)
c <- seq(from = 0, to = 10, by = 2)
v = 2L
w <<- 2 + 5i
a <- charToRaw("Hello")
\end{lstlisting}

The two \texttt{seq()} calls generate regularly spaced sequences. The suffix \texttt{L} in \texttt{2L} is used by the handout to create an integer value. The complex-number example contains a real and imaginary component. \texttt{charToRaw()} converts the characters of \texttt{"Hello"} to raw byte values, displayed in hexadecimal form in the handout.

\subsection{Inspecting an object's class}

The handout checks object classes with \texttt{class()}:

\begin{lstlisting}[style=dsr]
print(class(Gender))
print(class(height))
print(class(f))
print(class(x))
print(class(b))
print(class(v))
print(class(w))
\end{lstlisting}

\begin{keyidea}
The same variable syntax can hold different kinds of values. When an operation depends on the kind of value stored, inspect the object using \texttt{class()} rather than assuming its type from the variable name.
\end{keyidea}

\section{Arithmetic operations}

The handout introduces the following operators:

\begin{dstable}
\begin{tabularx}{\linewidth}{@{}p{2.2cm}X@{}}
\toprule
\rowcolor{DSPaleBlue!92}
Operator & Operation \\
\midrule
\texttt{+} & Addition \\
\texttt{-} & Subtraction \\
\texttt{*} & Multiplication \\
\texttt{**} or \texttt{\^} & Exponentiation \\
\texttt{/} & Division \\
\texttt{\%\%} & Modulus / remainder \\
\texttt{\%/\%} & Integer division \\
\bottomrule
\end{tabularx}
\end{dstable}

The lab reuses earlier variables to make each operator concrete:

\begin{lstlisting}[style=dsr]
print(f + 3)
print(height - x)
print(Weight * 2)
print(b ** 2)
print(c ^ 5)

m = height / 100
print(Weight / (m ** 2))
BMI = Weight / (m ** 2)

print(b %% 2)
print(c %/% 2)
\end{lstlisting}

The BMI sequence is especially useful because it combines assignment, division, exponentiation, and reuse of a previously calculated value.

\section{Reading user input}

The handout uses \texttt{readline()} for keyboard input and then converts the age value to numeric form:

\begin{lstlisting}[style=dsr]
name <- readline(prompt = "Enter name: ")
age <- readline(prompt = "Enter age: ")

# convert character into numeric
age <- as.numeric(age)

print(paste("Hi,", name,
            "this year you are", age, "years old."))
\end{lstlisting}

This example combines input, type conversion, text construction, and output.

\section{Built-in help and demonstrations}

The final Lab 1a exercises introduce R's help and demonstration commands:

\begin{lstlisting}[style=dsr]
?paste
demo(graphics)
demo(image)
\end{lstlisting}

These commands are part of the learning workflow: inspect function help, then run built-in demonstrations to see how features behave.

\begin{quickcheck}
\begin{enumerate}
  \item Why does the handout install R before RStudio?
  \item What is the practical difference between testing one expression in the Console and saving several statements in an R Script?
  \item What does \texttt{ls(pattern = "var")} do differently from \texttt{ls()}?
  \item Which operators in the handout return a remainder and an integer quotient?
  \item Why does the user-input example call \texttt{as.numeric(age)}?
\end{enumerate}
\end{quickcheck}

\begin{chapterrecap}
\begin{itemize}
  \item Verify the environment with a simple Console expression before starting larger exercises.
  \item Use scripts for saved, rerunnable programs and the Console for immediate evaluation.
  \item R supports leftward and rightward assignment forms in the supplied exercises.
  \item \texttt{class()} inspects an object's class; \texttt{ls()} inspects workspace object names.
  \item The arithmetic examples lead naturally into a small BMI calculation.
  \item \texttt{readline()}, \texttt{as.numeric()}, and \texttt{paste()} form a basic input-process-output pattern.
\end{itemize}
\end{chapterrecap}
